\section{Uniforme Continua}

La distribuzione \emph{Uniforme continua} assegna la stessa densità di probabilità
a tutti i valori appartenenti all’intervallo $[a,b]$.
In altre parole, ogni sottointervallo di uguale lunghezza ha la stessa probabilità
di realizzarsi.

\[
\mathcal{U}(a,b)
\]

\begin{figure}[H]
\centering

\begin{minipage}{0.45\textwidth}
\centering
\begin{tikzpicture}
\begin{axis}[
    axis lines = middle,
    ylabel = {$f(x)$},
    xlabel = {$x$},
    width = 0.90\textwidth,
    height = 0.90\textwidth,
    ymin = 0, ymax = 0.12,
    xmin = -1, xmax = 26,
    legend style={at={(axis cs: 18,0.11)}, anchor=north west, font=\tiny},
    legend cell align=left,
]

% --- densità dist 1: U(0,20), h=0.05 ---
\addplot[
    black,
    thick
] coordinates {
(-1,0) (0,0) (0,0.05) (20,0.05) (20,0) (26,0)
};
%\addlegendentry{$a=0,\, b=20$}

% --- densità dist 2: U(5,15), h=0.10 ---
\addplot[
    gray!60,
    thick
] coordinates {
(-1,0) (5,0) (5,0.10) (15,0.10) (15,0) (26,0)
};
%\addlegendentry{$a=5,\, b=15$}

\end{axis}
\end{tikzpicture}
\end{minipage}
\hfill
\begin{minipage}{0.45\textwidth}
\centering

\begin{tikzpicture}
\begin{axis}[
    axis lines = middle,
    ylabel = {$F(x)$},
    xlabel = {$x$},
    width = 0.90\textwidth,
    height = 0.90\textwidth,
    ymin = 0, ymax = 1.1,
    xmin = -1, xmax = 26,
    legend style={at={(axis cs: 18,0.25)}, anchor=north west, font=\tiny},
    legend cell align=left,
]

% --- CDF dist 1: U(0,20) ---
\addplot[
    black,
    thick
] coordinates {
(-1,0) (0,0)
(20,1) (26,1)
};
\addlegendentry{$a=0,\, b=20$}

% --- CDF dist 2: U(5,15) ---
\addplot[
    gray!60,
    thick
] coordinates {
(-1,0) (5,0)
(15,1) (26,1)
};
\addlegendentry{$a=5,\, b=15$}

\end{axis}
\end{tikzpicture}

\end{minipage}

\end{figure}


\bigskip
\begin{table}[H]
\centering
\renewcommand{\arraystretch}{2}
\begin{tabular}{@{}ll@{}}
\toprule
\textbf{Supporto} 
& $[a,b]\subset\mathbb{R}$ \\

\midrule

\textbf{Funzione di densità} 
& $
f_X(x)=
\begin{cases}
\dfrac{1}{b-a}, & a \le x \le b,\\[6pt]
0, & \text{altrimenti}
\end{cases}
$ \\

\midrule

\textbf{Funzione di ripartizione}
& $
F_X(x)=
\begin{cases}
0, & x < a,\\[6pt]
\dfrac{x-a}{b-a}, & a \le x \le b,\\[6pt]
1, & x > b
\end{cases}
$ \\

\midrule

\textbf{Valore atteso} 
& $\mathbb{E}[X] = \dfrac{a+b}{2}$ \\

\midrule

\textbf{Varianza}
& $Var(X) = \dfrac{(b-a)^2}{12}$ \\

\midrule

\textbf{Funzione caratteristica}
& $
\varphi_X(u)
=
\dfrac{e^{iub}-e^{iua}}{iu(b-a)}
$ \\

\bottomrule
\end{tabular}
\end{table}
